% from ega4-2-fr.tex
\oldpage[IV-2]{56}

\item[(iii)] Si en outre $X$ et $Y$ sont localement de type fini sur $k$,
toute composante irréductible de $X\times_k Y$ est de dimension
$\dim(X)+\dim(Y)$.
\end{enumerate}
\end{proposition}

\begin{proof}
L'assertion (iii) résulte de (4.2.1.2), compte tenu de (i) et (ii). Pour
démontrer (ii), on peut évidemment se borner au cas où
$X=\operatorname{Spec}(A)$, $Y=\operatorname{Spec}(B)$ sont affines, $A$ et
$B$ étant donc des anneaux intègres de corps des fractions respectifs
$K=R(X)$ et $L=R(Y)$~; l'assertion (i) montre que tout idéal premier minimal
$\mathfrak q$ de $A\otimes_k B$ est la trace sur $A\otimes_k B$ d'un idéal
premier minimal $\mathfrak p$ de $K\otimes_k L$~; comme $K\otimes_k L$ est un
anneau de fractions de $A\otimes_k B$, l'isomorphie des anneaux
$(A\otimes_k B)_{\mathfrak q}$ et $(K\otimes_k L)_{\mathfrak p}$ résulte de
$(0_{\mathrm I},1.2.5)$. Enfin, si par exemple $R(Y)$ est séparable sur $k$,
on sait que l'anneau $R(X)\otimes_k R(Y)$ est réduit (Bourbaki,
\emph{Alg.}, chap.~VIII, \S~7, n\textsuperscript{o}~3, th.~1)~; il en est donc
de même des anneaux locaux des points maximaux de $X\times_k Y$. On en déduit
que $X\times_k Y$ est réduit (3.2.1).
\end{proof}

\begin{proposition}[4.2.5]
\label{IV.2.4.2.5-fr}
Soient $k$ un corps, $X$ et $Y$ des $k$-préschémas localement noethériens,
$\mathcal F$ (resp. $\mathcal G$) un $\mathcal O_X$-Module (resp. un
$\mathcal O_Y$-Module) quasi-cohérent. Soit $(Z'_\lambda)$ (resp.
$(Z''_\mu)$) la famille des cycles premiers associés à $\mathcal F$ (resp.
$\mathcal G$), et désignons encore par $Z'_\lambda$ (resp. $Z''_\mu$) le
sous-préschéma réduit de $X$ (resp. $Y$) ayant $Z'_\lambda$ (resp.
$Z''_\mu$) comme espace sous-jacent. Alors, si $Z'_\lambda\times_k Z''_\mu$
est localement noethérien, les composantes irréductibles $Z_{\lambda\mu\nu}$
de $Z'_\lambda\times_k Z''_\mu$ dominent $Z'_\lambda$ et $Z''_\mu$, et
$(Z_{\lambda\mu\nu})$ est la famille des cycles premiers distincts associés à
$\mathcal F\otimes_k\mathcal G$.
\end{proposition}

\begin{proof}
Il suffit d'appliquer (4.2.4) au produit
$\operatorname{Spec}(k(x))\times_k\operatorname{Spec}(k(y))$.
\end{proof}

En particulier~:

\begin{corollary}[4.2.6]
\label{IV.2.4.2.6-fr}
Soient $k$ un corps, $X$, $Y$ deux $k$-préschémas localement noethériens tels
que $X\times_k Y$ soit localement noethérien. Soit $(Z'_\lambda)$ (resp.
$(Z''_\mu)$) la famille des sous-préschémas réduits de $X$ (resp. $Y$) ayant
pour espaces sous-jacents les composantes irréductibles de $X$ (resp. $Y$).
Alors les composantes irréductibles $Z_{\lambda\mu\nu}$ de
$Z'_\lambda\times_k Z''_\mu$ dominent $Z'_\lambda$ et $Z''_\mu$, et
$(Z_{\lambda\mu\nu})$ est la famille des composantes irréductibles de
$X\times_k Y$.
\end{corollary}

\begin{proof}
En effet, on peut se borner au cas où $X$ et $Y$ sont réduits (I, 5.1.8)~; les
composantes irréductibles de $X$ (resp. $Y$) sont alors les cycles premiers
associés à $\mathcal O_X$ (resp. $\mathcal O_Y$) (3.2.1). Appliquons (4.2.5)
à $\mathcal F=\mathcal O_X$ et $\mathcal G=\mathcal O_Y$, en notant que par
définition $\mathcal O_X\otimes_k\mathcal O_Y=\mathcal O_{X\times_k Y}$~; le
corollaire résulte de ce que $\mathcal O_X\otimes_k\mathcal O_Y$ n'a pas de
cycles premiers associés immergés puisqu'il en est ainsi de $\mathcal O_X$ et
$\mathcal O_Y$ par hypothèse (4.2.3).
\end{proof}

Appliquons les résultats qui précèdent au cas où $Y$ est le spectre d'une
extension $K$ de $k$~:

\begin{proposition}[4.2.7]
\label{IV.2.4.2.7-fr}
Soient $k$ un corps, $X$ un $k$-préschéma, $K$ une extension de $k$ telle que
$X\otimes_k K$ soit localement noethérien, $\mathcal F$ un
$\mathcal O_X$-Module quasi-cohérent, $x'$ un point de $X\otimes_k K$, $x$
son image dans $X$.
\begin{enumerate}
\item[(i)] Soit $(Z_\lambda)$ la famille des sous-préschémas réduits de $X$
ayant pour espaces sous-jacents les cycles premiers associés à $\mathcal F$~;
alors les composantes irréductibles $Z_{\lambda\mu}$ des
$Z_\lambda\otimes_k K$ sont les cycles premiers associés à
$\mathcal F\otimes_k K$, et $Z_{\lambda\mu}$ domine $Z_\lambda$~; en outre,
pour qu'un $Z_{\lambda\mu}$ soit immergé, il faut et il suffit que $Z_\lambda$
le soit.
\item[(ii)] Pour que $x$ appartienne à un cycle premier associé immergé de
$\mathcal F$, il faut et il suffit que $x'$ appartienne à un cycle premier
associé immergé de $\mathcal F\otimes_k K$~; pour que $\mathcal F$ soit sans
cycle premier associé immergé, il faut et il suffit qu'il en soit ainsi de
$\mathcal F\otimes_k K$.

